\documentclass[10pt,a4paper]{amsart}
\usepackage[margin=19mm]{geometry}
\usepackage{amsmath,amssymb,mathtools}
\usepackage[scr=rsfs,cal=euler]{mathalfa}
\usepackage{xfrac}
\usepackage[unicode,hidelinks,bookmarks=false]{hyperref}
\newcommand{\ie}{i.e.\ }
\newcommand{\cf}{cf.\ }
\newcommand{\resp}{respectively\ }
\newcommand{\Subsection}{\subsection}
\providecommand{\nobreakdash}{\nobreak\hbox{-}}
\setlength{\emergencystretch}{2em}
\multlinegap0pt
\usepackage[shortlabels]{enumitem}
\usepackage{amssymb}
\usepackage{mathtools}
\usepackage{dsfont}
\usepackage{float}
\usepackage{xcolor}
\newtheorem{proposition}{Proposition}
\title{Learning on the Temporal Tangent Bundle for Physics-Informed Neural Networks}
\author{Adetola Jamal, Mamlankou Charbel, Hou\'edanou Koffi Wilfrid, D\`egla Aymard Guy}
\date{Source: arXiv:2604.11829v1 (CC BY 4.0)}
\begin{document}
\maketitle

\noindent\emph{Source and licence.} Learning on the Temporal Tangent Bundle for Physics-Informed Neural Networks. Version arXiv:2604.11829v1 (CC BY 4.0), distributed by arXiv under the Creative Commons Attribution 4.0 International licence, http://creativecommons.org/licenses/by/4.0/, as recorded in the arXiv licence metadata for this exact version and shown on the abstract page https://arxiv.org/abs/2604.11829v1. Full licence notice: \texttt{licenses/arxiv-2604.11829-CC-BY-4.0.txt}.\par\smallskip

\noindent\textit{Editorial scope.} Counted unit: the article's proposition prop:condition\_number (source0.tex lines 1132--1148). The article's complete original proof of it is delivered with it (source0.tex lines 1150--1194, one \texttt{proof} environment of 45 lines). No mathematics is written, repaired or re-derived: the statement and the proof are the article's own lines, in the article's own notation, and no retained line is substituted or re-flowed.  Retained uncounted as the source-local context the proof needs: lines 307--312.  Nothing was omitted from any retained span, so the package carries no omission record.  No retained line contains a citation, so the wrapper carries no bibliography and no bibliography entry.  No retained line contains a figure environment, an image-inclusion command or drawing code, so no image is delivered and the package declares no asset dependency.  Editorial formatting only, in addition: an AMS \texttt{amsart} preamble that restates the article's own declarations and macros used by the retained lines, namely the environments \texttt{proposition} on separate counters, together with the standard packages those lines need.  The retained environments print in this document as Proposition 1 in order of appearance, whereas the article prints the same items with the numbering of its own full document.  The complete native source of the article as distributed in the arXiv e-print for this exact version is bundled unchanged as \texttt{sources/198366/source0.tex}.
\begin{proposition}[Continuity and Lipschitz property]\label{prop:continuity}
The operator $\mathcal{I}_{u_0}: L^2(0,T; L^2(\Omega)) \to C([0,T]; L^2(\Omega))$ is continuous and Lipschitz continuous with constant $L = \sqrt{T}$. Specifically, for any $v_1, v_2 \in L^2(0,T; L^2(\Omega))$,
\begin{equation}
    \sup_{t \in [0,T]} \|\mathcal{I}_{u_0}[v_1](\cdot,t) - \mathcal{I}_{u_0}[v_2](\cdot,t)\|_{L^2(\Omega)} \leq \sqrt{T} \|v_1 - v_2\|_{L^2(0,T; L^2(\Omega))}.
\end{equation}
\end{proposition}

\begin{proposition}[Condition number and round-off error accumulation]\label{prop:condition_number}
Define the \textbf{condition number} of the Volterra operator as
\begin{equation}
    \kappa(\mathcal{I}_{u_0}) := \|\mathcal{I}_{u_0}\|_{\mathrm{op}} \cdot \|\mathcal{I}_{u_0}^{-1}\|_{\mathrm{op}},
\end{equation}
where $\|\cdot\|_{\mathrm{op}}$ denotes the operator norm in $L^2(\Omega \times (0,T])$. Then:

\begin{enumerate}[label=(\alph*)]
    \item The condition number satisfies $\kappa(\mathcal{I}_{u_0}) = \sqrt{T}$, which is mild and grows only as the square root of the integration time.
    
    \item Suppose round-off errors of magnitude $\varepsilon_{\mathrm{mach}}$ (machine precision) are introduced at each quadrature point evaluation. The accumulated error after $N = MT$ time steps satisfies
    \begin{equation}\label{eq:roundoff_accumulation}
        \|\text{accumulated error}\|_{L^2(\Omega)} \leq C_{\mathrm{acc}} \sqrt{N} \, \varepsilon_{\mathrm{mach}} = C_{\mathrm{acc}} \sqrt{MT} \, \varepsilon_{\mathrm{mach}},
    \end{equation}
    demonstrating sub-linear growth in the number of steps.
\end{enumerate}
\end{proposition}

\begin{proof}
\textbf{Part (a) - Condition number:}

\noindent From Proposition~\ref{prop:continuity}, the forward operator $\mathcal{I}_{u_0}: L^2(\Omega \times (0,T]) \to L^\infty(0,T; L^2(\Omega))$ satisfies
\begin{equation}
    \|\mathcal{I}_{u_0}[v]\|_{L^\infty(0,T; L^2(\Omega))} \leq \sqrt{T} \|v\|_{L^2(\Omega \times (0,T])},
\end{equation}
so $\|\mathcal{I}_{u_0}\|_{\mathrm{op}} = \sqrt{T}$.

\noindent The inverse operator $\mathcal{I}_{u_0}^{-1}$ is the differentiation operator $\mathcal{D}[\cdot] = \partial_t[\cdot]$ (restricted to the subspace of functions satisfying $u(0) = u_0$). In the $L^2$ sense, differentiation is an isometry modulo boundary terms:
\begin{equation}
    \|\mathcal{D}[u]\|_{L^2(\Omega \times (0,T])}^2 = \int_0^T \|\partial_t u(\cdot,t)\|_{L^2(\Omega)}^2 \, \mathrm{d}t.
\end{equation}
For functions $u \in \mathcal{V} = C^1([0,T]; H^k(\Omega))$ with $u(0) = u_0$, we have $\|\mathcal{D}\|_{\mathrm{op}} = 1$. Therefore,
\begin{equation}
    \kappa(\mathcal{I}_{u_0}) = \|\mathcal{I}_{u_0}\|_{\mathrm{op}} \cdot \|\mathcal{D}\|_{\mathrm{op}} = \sqrt{T} \cdot 1 = \sqrt{T}.
\end{equation}

\medskip
\noindent\textbf{Part (b) - Round-off error accumulation:}

\noindent Consider the discrete trapezoidal rule where at each quadrature point $s_m = m h$ (with $h = T/(MT) = 1/M$), the velocity evaluation $v_\theta(x, s_m)$ incurs a round-off error $\varepsilon_m(x)$ with $\|\varepsilon_m\|_{L^2(\Omega)} \leq \varepsilon_{\mathrm{mach}}$. The accumulated error in the reconstructed state at time $t = Nh$ (after $N$ steps) is
\begin{equation}
    \text{accumulated error} = \sum_{m=0}^{N-1} \frac{h}{2} \left( \varepsilon_m + \varepsilon_{m+1} \right).
\end{equation}
In the worst case (deterministic alignment of errors), the $L^2$ norm satisfies
\begin{equation}
    \|\text{accumulated error}\|_{L^2(\Omega)} \leq \sum_{m=0}^{N-1} \frac{h}{2} \left( \|\varepsilon_m\|_{L^2} + \|\varepsilon_{m+1}\|_{L^2} \right) \leq \sum_{m=0}^{N-1} h \varepsilon_{\mathrm{mach}} = Nh \varepsilon_{\mathrm{mach}} = t \varepsilon_{\mathrm{mach}}.
\end{equation}
This gives a linear accumulation in time, which is typical for numerical integration.

\noindent However, if we assume the round-off errors $\{\varepsilon_m\}$ are independent random variables with zero mean and variance $\sigma^2 = \varepsilon_{\mathrm{mach}}^2$, then by the central limit theorem, the variance of the accumulated error is
\begin{equation}
    \mathbb{E}\left[ \|\text{accumulated error}\|_{L^2}^2 \right] = \sum_{m=0}^{N-1} \left( \frac{h}{2} \right)^2 \left( \sigma^2 + \sigma^2 \right) = N \frac{h^2}{2} \sigma^2 = \frac{Nh^2}{2} \varepsilon_{\mathrm{mach}}^2.
\end{equation}
Taking the square root yields the root-mean-square error
\begin{equation}
    \sqrt{\mathbb{E}\left[ \|\text{accumulated error}\|_{L^2}^2 \right]} = \frac{h\varepsilon_{\mathrm{mach}}}{\sqrt{2}} \sqrt{N} = \frac{\varepsilon_{\mathrm{mach}}}{\sqrt{2M}} \sqrt{N}.
\end{equation}
Since $N = MT$ and $h = 1/M$, this simplifies to
\begin{equation}
    \sqrt{\mathbb{E}\left[ \|\text{accumulated error}\|_{L^2}^2 \right]} = \frac{\varepsilon_{\mathrm{mach}}}{\sqrt{2M}} \sqrt{MT} = \varepsilon_{\mathrm{mach}} \sqrt{\frac{T}{2M}} \sqrt{M} = \varepsilon_{\mathrm{mach}} \sqrt{\frac{T}{2}} \approx 0.707 \sqrt{T} \, \varepsilon_{\mathrm{mach}}.
\end{equation}
This is independent of $M$ and grows only as $\sqrt{T}$, matching the condition number. For practical purposes, we can write this as $C_{\mathrm{acc}} \sqrt{MT} \varepsilon_{\mathrm{mach}}$ where $C_{\mathrm{acc}} \sim 1/\sqrt{M}$ accounts for the averaging effect of the trapezoidal rule, yielding the sub-linear growth stated in~\eqref{eq:roundoff_accumulation}.
\end{proof}

\end{document}
